\ifdefined\CRevisionRound\else\def\CRevisionRound{2}\fi
\documentclass[10pt]{article}
\pdfvariable suppressoptionalinfo 767
\usepackage[a4paper,margin=18mm]{geometry}
\usepackage{amsmath,amssymb,mathtools,microtype}
\usepackage{fontspec}
\setmainfont{TeX Gyre Pagella}
\setsansfont{TeX Gyre Heros}
\setmonofont{Latin Modern Mono}
\usepackage[hidelinks]{hyperref}
\newcommand{\Orb}{\mathcal O}
\newcommand{\RR}{\mathbb R}
\newcommand{\ZZ}{\mathbb Z}
\newtheorem{theorem}{Theorem}
\title{The \(U(3)\) Gelfand--Tsetlin System:\\Exact Fibers, Linear Dynamics, and Unshifted Quantization}
\author{HCS-C365 theorem package}
\date{4 September 2026}
\begin{document}
\maketitle
\begin{abstract}
We close the classical and quantum Gelfand--Tsetlin descriptions for a Hermitian $U(3)$ coadjoint orbit under explicit conventions.  An arrow completion proves the exact interlacing image; nested-group flows give every regular Lagrangian torus; character duality classifies every linear Hamiltonian closure and least period.  With integral highest weight and no half-form correction, branching gives the complete basis with unshifted GT labels and the Weyl count.  Repeated weights, facets, and nonintegral spectra are separated sharply.  Exact enumeration is a regression receipt, not the proof.
\end{abstract}

\section{Frozen orbit and exact theorem}
Let $\lambda_1\geq\lambda_2\geq\lambda_3$, and let $\Orb_\lambda$ be the orbit of Hermitian matrices with this spectrum.  For skew-Hermitian $X,Y$, freeze
\[
 \omega_A([X,A],[Y,A])=-i\operatorname{Tr}(A[X,Y]),\qquad
 f_X(A)=-i\operatorname{Tr}(AX),\qquad \iota_{X_{f_X}}\omega=df_X.
\]
Thus $X_{f_X}(A)=[X,A]$.  Write $B=A^{(2)}$ for the northwest $2\times2$ minor, let $y_1\geq y_2$ be its eigenvalues, and put $x=A_{11}$.

\begin{theorem}\label{thm:main}
The map $F=(y_1,y_2,x)$ has image
\[
 \Delta_\lambda=\{\lambda_1\geq y_1\geq\lambda_2\geq y_2\geq\lambda_3,
 \quad y_1\geq x\geq y_2\}.
\]
If $\lambda_1>\lambda_2>\lambda_3$, every strict-interlacing fiber is exactly an effective Lagrangian $T^3$-orbit, with each circle normalized to period $2\pi$.

For $H_\nu=\nu_1y_1+\nu_2y_2+\nu_3x$, its regular flow is
$\theta(t)=\theta_0+t\nu\pmod {2\pi}$.  Let $K_\nu$ be the subgroup annihilated by every $k\in\ZZ^3$ satisfying $k\cdot\nu=0$.  The orbit closure is the coset $\theta_0+K_\nu$, and $\theta_0+K_\nu=K_\nu$ exactly when $\theta_0\in K_\nu$, with
$\dim K_\nu=\operatorname{rank}_{\mathbb Q}\{\nu_1,\nu_2,\nu_3\}$.  It is fixed for $\nu=0$.  For $\nu\ne0$, it is periodic precisely when $\nu=\alpha m$, where $\alpha>0$ and $m\in\ZZ^3$ is primitive; its least period is $2\pi/\alpha$.

If $\lambda\in\ZZ^3$ is dominant, take prequantum curvature $-i\omega$ and no half-form correction.  Borel--Weil quantization has a multiplicity-free GT branching basis indexed, without a $\rho$-shift, by
\[
 \lambda_1\geq m_{21}\geq\lambda_2\geq m_{22}\geq\lambda_3,
 \qquad m_{21}\geq m_{11}\geq m_{22},
\]
with joint branching labels $(m_{21},m_{22},m_{11})$, and
\[
 \dim Q(\Orb_\lambda)=
 \frac{(\lambda_1-\lambda_2+1)(\lambda_2-\lambda_3+1)
 (\lambda_1-\lambda_3+2)}2.
\]
\end{theorem}

\section{Interlacing and constructive surjectivity}
Cauchy interlacing gives all six inequalities.  Conversely, at any point of $\Delta_\lambda$, set
\[
 B_{11}=x,\quad B_{12}=B_{21}=r,\quad B_{22}=y_1+y_2-x,
 \qquad r^2=(y_1-x)(x-y_2).
\]
This has eigenvalues $y_1,y_2$.  In a basis diagonalizing $B$, consider
\[
 A'_{11}=y_1,\ A'_{22}=y_2,\ A'_{12}=0,\ A'_{13}=z_1,\ A'_{23}=z_2,
 \quad A'_{33}=c=\lambda_1+\lambda_2+\lambda_3-y_1-y_2,
\]
with the lower-triangular entries fixed by Hermitian symmetry.
If $P(t)=\prod_{j=1}^3(t-\lambda_j)$, choose
\[
 |z_1|^2=-\frac{P(y_1)}{y_1-y_2},\qquad
 |z_2|^2=\frac{P(y_2)}{y_1-y_2}.
\]
Interlacing makes both quantities nonnegative.  The monic cubic $\det(tI-A')$ has the correct trace and agrees with $P$ at $y_1,y_2$, hence equals $P$.  Conjugating the first two coordinates back retains $A_{11}=x$.  Endpoint coincidences follow by setting the vanishing arrow coordinate to zero, or by continuity.  This proves the exact image.

\ifnum\CRevisionRound>0
\section{Thimm flows and complete linear dynamics}
Where $B$ is simple, $dy_j$ is its spectral projector in the northwest block.  The Hamiltonian flow is conjugation by $e^{itP_j}$, while the $x$-flow is conjugation by $e^{itE_{11}}$.  Each has period $2\pi$.  Spectral functions on the same $U(2)$ block commute; invariance along the nested $U(1)\subset U(2)$ action gives commutation with $x$.  Strict interlacing makes the three infinitesimal actions independent.

For completeness, a regular fiber has no hidden component.  In the arrow construction, strictness makes $r,z_1,z_2$ nonzero.  After fixing an eigenvector gauge, their three independent phase choices form one $T^3$ and exhaust every matrix having the frozen $x,y_1,y_2$; changing the gauge merely relabels these phases.  Thus the fiber is precisely the compact three-dimensional orbit.  It is half-dimensional and isotropic, hence Lagrangian.

The linear-flow assertion now follows on $\RR^3/(2\pi\ZZ)^3$.  The closure through the origin is the compact subgroup $K_\nu$ annihilated by
\(
 \{k\in\ZZ^3:k\cdot\nu=0\}.
\)
Character duality gives $\dim K_\nu=\operatorname{rank}_{\mathbb Q}\{\nu_j\}$.  Translating by $\theta_0$ gives the general closure $\theta_0+K_\nu$, which equals $K_\nu$ if and only if $\theta_0\in K_\nu$.  Thus a nonzero basepoint can still give the subgroup: for rational rank three $K_\nu=T^3$, while for $\nu=(1,2,3)$ the point $(\pi/2,\pi,3\pi/2)$ lies in $K_\nu$.  By contrast $(0,0,\pi)$ does not.  Rank one gives $\nu=\alpha m$ with $m$ primitive, and $t\nu\in2\pi\ZZ^3$ first at $t=2\pi/\alpha$.

\section{Boundary atlas}
For a regular spectrum, every polytope facet is singular: at least one circle stabilizer occurs, and no universal topology is claimed at facet intersections.  If $\lambda_1=\lambda_2>\lambda_3$, then $y_1=\lambda_1$; if $\lambda_1>\lambda_2=\lambda_3$, then $y_2=\lambda_3$.  In either case the orbit is four-dimensional, the image is a triangle, and only its regular interior carries the effective $T^2$ statement.  A scalar spectrum gives a point.  A nonintegral ordered spectrum retains the classical theorem but is outside the quantization assertion.
\fi

\ifnum\CRevisionRound>1
\section{Unshifted branching quantization}
With the frozen curvature and no metaplectic correction, Borel--Weil gives the irreducible $U(3)$-module of highest weight $\lambda$, including negative $\lambda_3$ through a determinant twist.  Restriction along
$U(3)\supset U(2)\supset U(1)$ is multiplicity free.  A $U(2)$ weight $(m_{21},m_{22})$ occurs exactly under the first interlacing row, and $m_{11}$ occurs exactly under the second.  Thus there is one branching vector per displayed pattern, with unshifted joint labels; no spectrum of unspecified Casimir or quantum action operators is claimed.

Put $a=\lambda_1-\lambda_2$ and $b=\lambda_2-\lambda_3$.  Translation by $\lambda_3$ preserves the count, and direct summation yields
\[
 \sum_{u=0}^{a}\sum_{v=0}^{b}(b+u-v+1)
 =\frac{(a+1)(b+1)(a+b+2)}2,
\]
which is the stated dimension.  This also closes repeated weights: $a=0$ or $b=0$ gives the triangular $T^2$ count, while $a=b=0$ gives one vector.

\section{Exact receipt, sources, and Route-A boundary}
For every $0\leq a,b\leq12$, the artifact exhausts all patterns for $\lambda=(a+b,b,0)$: 169 weights and 74,529 patterns, committed by per-weight canonical hashes.  Eight formal frequency triples test ranks zero through three and exact least-period tokens; three basepoint rows lock the exact coset-equality criterion; six rows lock the boundary atlas.  An independent implementation reconstructs all rows, while SymPy checks the arrow cubic, count, KKS identity, moment pairing, projector period, and the basepoint witnesses.  The hostile suite repairs hashes after semantic and type mutations before requiring rejection.  These checks are finite regression evidence only.

Guillemin and Sternberg's flag-system paper (J. Funct. Anal. 52 (1983), DOI
\href{https://doi.org/10.1016/0022-1236(83)90092-7}{10.1016/0022-1236(83)90092-7}) and their multiplicity paper (Invent. Math. 67 (1982), DOI
\href{https://doi.org/10.1007/BF01398934}{10.1007/BF01398934}) provide primary lineage.  We claim no priority and have reconstructed the $U(3)$ statements under the conventions above.

The integral branching lattice supports only a weak arithmetic relation; the torus flow supports only weak dynamics.  Principal-minor polynomials are not target determinants.  There is no target Euler factor, root number, automorphy, divisor, functional equation, target-zero match, or Hilbert--Pólya operator.  Route B is false.  Scope is \texttt{NO\_BAD\_EULER\_OR\_ROOT\_NUMBER}.
\fi

\vfill\noindent\footnotesize
\ifcase\CRevisionRound
round zero interlacing and arrow completion closure
\or round one thimm torus and closure rank
\else round two branching quantization and route a closure
\fi
\end{document}
